\subsection{Experimental Setup}

All experiments were run on a DigitalOcean droplet with two AMD virtual CPU cores at 2.00 GHz and 4 GiB of memory.
The operating system was Debian GNU/Linux 13.4 (trixie), using Linux kernel 6.12.43.
All models were trained and evaluated on CPUs without GPU acceleration.
All experiments were conducted using Python 3.12.11 with NumPy 2.4.4, pandas 3.0.3, scikit-learn 1.8.0, statsmodels 0.14.6, PyTorch 2.11.0, and Matplotlib 3.10.8.

For each imputation run, two interior observations were randomly held out for testing, while the first and last observations were retained in the training set.
Forecasting used rolling-origin evaluation with four origins per series, with five test observations per origin for global plastic production and two for Japanese plastic waste generation.
At each origin, models were trained on all preceding observations and evaluated on the subsequent test window.

For ridge regression, Gaussian process regression, and support vector regression, predictors and targets were standardized using parameters fitted separately on the training portion of each fold.
ARIMA and exponential smoothing models were fitted on the original scale.
For feature-importance analysis, isolation forest was fitted to the predictor variables before fitting the random forest.

During hyperparameter tuning, forecasting models were evaluated using time-series cross-validation with \(K = 4\) validation folds.
The chronologically ordered training data were split into \(K + 1\) blocks, with the first block serving as the initial training window.
Each subsequent block served once as the validation window, with the model trained on all preceding blocks.

Imputation models were evaluated using repeated random holdout validation with four masks.
Each mask reserved 20\% of the training samples, with at least two samples reserved for validation, and the remaining samples were used for fitting.
